\documentclass[aspectratio=169,10pt]{beamer}
\usetheme[numbering=fraction,progressbar=frametitle]{metropolis}
\usepackage{fontspec}
\setsansfont{Latin Modern Sans}
\usepackage{amsmath,amssymb,mathtools,graphicx,booktabs}
\usepackage{polyglossia}
\setdefaultlanguage{english}
\setbeamertemplate{navigation symbols}{}
\setbeamerfont{frametitle}{size=\large}
\setbeamerfont{normal text}{size=\normalsize}
\setlength{\parskip}{.45em}
\hypersetup{pdfauthor={},pdfsubject={Lecture slides},pdfkeywords={},pdftitle={Antiderivatives and differential equations}}
\providecommand{\N}{\mathbb{N}}
\providecommand{\Z}{\mathbb{Z}}
\providecommand{\D}{\mathbb{D}}
\providecommand{\Q}{\mathbb{Q}}
\providecommand{\R}{\mathbb{R}}
\providecommand{\C}{\mathbb{C}}
\providecommand{\vect}[1]{\overrightarrow{#1}}
\providecommand{\vi}{\vec{\imath}}
\providecommand{\vj}{\vec{\jmath}}
\providecommand{\vk}{\vec{k}}
\providecommand{\norm}[1]{\left\lVert #1\right\rVert}
\providecommand{\vc}[2]{\begin{pmatrix}#1\\#2\end{pmatrix}}
\providecommand{\deter}[2]{\det\left(#1,#2\right)}
\providecommand{\abs}[1]{\left\lvert #1\right\rvert}
\providecommand{\dd}{\,\mathrm{d}}
\providecommand{\e}{\mathrm{e}}
\providecommand{\Card}{\operatorname{Card}}
\providecommand{\E}{\operatorname{E}}
\providecommand{\V}{\operatorname{V}}
\providecommand{\degre}{^{\circ}}
\providecommand{\un}[1]{\,\mathrm{#1}}
\providecommand{\sci}[2]{#1\times10^{#2}}
\providecommand{\pgcd}{\operatorname{PGCD}}
\providecommand{\ppcm}{\operatorname{PPCM}}
\title{Antiderivatives and differential equations}
\subtitle{Lecture and worked examples}
\date{}
\author{}
\begin{document}
\begin{frame}[plain]
\vfill
{\usebeamerfont{title}\inserttitle\par}
\medskip
{\usebeamerfont{subtitle}\insertsubtitle\par}
\medskip
{\usebeamercolor[fg]{progress bar}\rule{\textwidth}{.5pt}}
\vfill
\end{frame}
\begin{frame}{Objectives}
Calculate antiderivatives; solve a differential equation with an initial condition; interpret an Euler approximation.
\end{frame}
\begin{frame}{Antiderivatives on an interval}
At this stage we accept that every continuous function on an interval has antiderivatives; the link with area will be established in chapter 8.
An antiderivative of \(f\) on an interval \(I\) is a differentiable function \(F\) such that \(F'=f\) on \(I\). If \(F\) is an antiderivative, all the antiderivatives on \(I\) are \(F+C\), \(C\in\R\).
\end{frame}
\begin{frame}{Antiderivatives on an interval}
If \(F'=G'=f\), then \((F-G)'=0\) on the interval; \(F-G\) is therefore constant. Conversely, \((F+C)'=f\). The interval hypothesis is essential: on a union of disjoint intervals, such as \(\mathbb{R}\setminus\{0\}\), a union of two intervals, the constant may differ from one interval to the other.
\end{frame}
\begin{frame}{Antiderivatives on an interval}
\[\begin{array}{c|c}
f(x)&F(x)\\\hline
x^n\ (n\in\Z,\ n\ne-1)&\dfrac{x^{n+1}}{n+1}\\
\dfrac{1}{x}&\ln|x|\quad\text{sur un intervalle évitant }0\\
\dfrac{1}{\sqrt{x}}&2\sqrt{x}\quad(x>0)\\
e^x&e^x\\
\cos x&\sin x\\
\sin x&-\cos x
\end{array}\]
The additive constant can always be added. The last two lines will be revisited after the study of trigonometry.
\end{frame}
\begin{frame}{Recognising a composition}
For a differentiable function \(u\), the antiderivatives of \(u'u^n\), \(u'e^u\), \(\dfrac{u'}{u}\) are respectively \(\dfrac{u^{n+1}}{n+1}\) (\(n\ne-1\)), \(e^u\), \(\ln|u|\), on the intervals where the expressions make sense. Linearity allows us to handle sums.
\end{frame}
\begin{frame}{Recognising a composition}
An antiderivative of \(6x(x^2+1)^2\) is \((x^2+1)^3\); an antiderivative of \(e^{3x}\) is \(\dfrac{e^{3x}}{3}\). To impose \(F(0)=2\) in the latter case, take \(F(x)=\dfrac{e^{3x}}{3}+\dfrac{5}{3}\).
\end{frame}
\begin{frame}{Differential equations}
A differential equation relates an unknown function to its derivatives. A solution must be verified on a specified interval.
\end{frame}
\begin{frame}{Differential equations}
\(y'=ay\) has solutions \(y(x)=Ce^{ax}\). For \(a\ne0\), \(y'=ay+b\) has solutions
\[y(x)=Ce^{ax}-\dfrac ba.\]
If \(a=0\), the equation becomes \(y'=b\) and \(y(x)=bx+C\).
If \(g\) is a known particular solution of \(y'=ay+f(x)\), the solutions are \(y=g+Ce^{ax}\).
\end{frame}
\begin{frame}{Differential equations}
If \(y\) is a solution, \(z(x)=e^{-ax}y(x)\) satisfies \(z'=e^{-ax}(y'-ay)=0\); \(z=C\). Conversely, the derivative of \(Ce^{ax}\) is \(aCe^{ax}\). For \(y'=ay+b\), subtracting the constant solution \(-\dfrac{b}{a}\) reduces it to the homogeneous equation. For a variable right-hand side, do the same with \(g\).
\end{frame}
\begin{frame}{Differential equations}
\(y'=-2y+6\) and \(y(0)=1\) give \(y=3+Ce^{-2x}\), then \(C=-2\). The solution \(3-2e^{-2x}\) increases towards the equilibrium \(3\) as \(x\to+\infty\).
\end{frame}
\begin{frame}{Euler's method}
For \(y'=ay+b\), with step \(h>0\), we approximate \(y(t_k)\) by \(Y_k\) according to \(t_{k+1}=t_k+h\), \(Y_{k+1}=Y_k+h(aY_k+b)\). This is an approximation using the tangent, not an equality for the exact solution. A smaller step may improve the approximation; check the error numerically on a known solution.
\end{frame}
\begin{frame}{Research activity}
\textbf{Cooling of a drink - 40 min.}

A drink initially at \(80\,^{\circ}\mathrm C\) is placed in a room kept at \(20\,^{\circ}\mathrm C\). For \(t\ge0\), let \(T(t)\) denote its temperature in degrees Celsius, where \(t\) is the time elapsed in minutes. We accept the model \(T' (t)=-0.1(T(t)-20)\) and \(T(0)=80\).
\end{frame}
\begin{frame}{Research activity (2)}
\begin{enumerate}
\setcounter{enumi}{0}\item Solve this initial value problem and give the expression of \(T(t)\).

\item Study the direction of variation of \(T\) and find its limit at \(+\infty\). Interpret the results.

\item Find the first time at which the temperature equals \(40\,^{\circ}\mathrm C\). Give the exact value, then an approximate value in minutes.

\item Apply two steps of Euler's method with a step of \(1\) minute. Compare the approximation obtained after two minutes with the exact value \(T(2)\).

\end{enumerate}
\end{frame}
\begin{frame}{Research activity (3)}
\begin{enumerate}
\setcounter{enumi}{4}\item Give a concrete situation in which the assumption of a constant room temperature would no longer be appropriate.

\end{enumerate}
\end{frame}
\end{document}
